\documentclass{article}
\begin{document}
\textbf{1. Les fonctions :}
\begin{itemize}
\item Notion de fonction et représentation graphique
\item Variations et tableau de variations
\item Utilisation de fonctions pour modéliser des situations économiques ou de gestion
\item Interprétation de graphiques et de courbes de fonctions
\end{itemize}

\textbf{2. Les suites :}
\begin{itemize}
\item Généralités sur les suites arithmétiques et géométriques
\item Calcul de termes d'une suite
\item Utilisation de suites pour modéliser des situations économiques ou de gestion
\end{itemize}

\textbf{3. Les probabilités :}
\begin{itemize}
\item Notion de probabilité
\item Calcul de probabilités simples et conditionnelles
\item Utilisation de probabilités pour résoudre des problèmes économiques ou de gestion
\end{itemize}

\textbf{4. Les statistiques :}
\begin{itemize}
\item Collecte et organisation des données
\item Calcul de caractéristiques de position (moyenne, médiane, quartiles)
\item Calcul de caractéristiques de dispersion (étendue, variance, écart-type)
\item Interprétation de résultats statistiques
\end{itemize}

\textbf{5. Les pourcentages :}
\begin{itemize}
\item Calcul de pourcentages
\item Calcul de taux d'évolution et de coefficient multiplicateur
\item Utilisation des pourcentages dans des situations économiques ou de gestion
\end{itemize}

\textbf{6. Les fonctions affines :}
\begin{itemize}
\item Équations de droites et représentation graphique
\item Droites parallèles et perpendiculaires
\item Utilisation des fonctions affines pour modéliser des situations économiques ou de gestion
\end{itemize}

\textbf{7. Les équations et inéquations :}
\begin{itemize}
\item Résolution d'équations du premier degré
\item Résolution d'équations du second degré
\item Résolution d'inéquations du premier degré
\end{itemize}

\textbf{8. Les fonctions exponentielles et logarithmes :}
\begin{itemize}
\item Généralités sur les fonctions exponentielles et logarithmes
\item Résolution d'équations exponentielles et logarithmiques
\item Utilisation des fonctions exponentielles et logarithmes dans des situations économiques ou de gestion
\end{itemize}
\end{document}